% Notas de aula — Capítulo 17: Ventos Regionais
% Curso: Meteorologia Prática (baseado em R. Stull, Practical Meteorology, CC BY-NC-SA 4.0)
% Caixas verdes = conteúdo literal do quadro; linhas verdes = encenação.
\documentclass[11pt]{article}
\usepackage{../comum/preambulo}
\graphicspath{{../figuras/cap17/}}

\title{Capítulo 17 --- Ventos Regionais\\
{\large Notas de aula (roteiro de quadro-negro)}}
\author{Curso de Meteorologia Prática}
\date{}

\begin{document}
\maketitle
\thispagestyle{fancy}

\noindent\textbf{Plano sugerido:} 2 aulas de 2\,h.
\emph{Aula 1:} \S\ref{sec:termicos}--\S\ref{sec:cat} (circulações
térmicas: brisas, vales e catabáticos).
\emph{Aula 2:} \S\ref{sec:ondas}--\S\ref{sec:hidra} (ondas de montanha,
föhn, dinâmica de água rasa e vendavais).

\section{O tempo feito em casa}\label{sec:termicos}

\begin{noquadro}[abertura]
\textbf{Cap.~17 --- Ventos Regionais}\\[2pt]
entre a escala sinótica (Caps.~12--13) e a turbulência (Cap.~18):\\
os ventos que a \emph{geografia local} fabrica\\[2pt]
duas famílias: \textbf{térmicos} (contraste de aquecimento) e
\textbf{dinâmicos} (montanha no caminho)
\end{noquadro}

Quando a forçante sinótica está fraca (dias de alta
\javimos{12}{anticiclones}), a geografia assume: mar/terra, encosta,
vale e serra criam seus próprios ventos. É a meteorologia que se
\emph{sente} morando em Santos, Campos do Jordão ou no vale do
Paraíba.

\subsection{Brisa marítima}

\quadro{Desenhar o circuito vertical costa--mar com o solenoide:
colunas de densidades diferentes $\Rightarrow$ o circuito gira.
Anunciar: ``é um teorema de circulação, não um truque''.}

\begin{formadif}[teorema de Bjerknes (T1; Caso~1)]
\[ \frac{d\Gamma}{dt} = -\oint\frac{dP}{\rho} \neq 0
\quad\text{quando isóbaras cruzam isopicnas.} \]
Terra $\sim$5\,K mais quente de dia \javimos{3}{Bowen} monta uma
brisa de 3--7\,m/s em 1--2\,h, freada pelo atrito. A frente de brisa
avança terra adentro (10--50\,km) como mini-frente fria
\javimos{12}{frentes} — e dispara a convecção da tarde
\javimos{14}{gatilhos}.
\end{formadif}

\begin{noquadro}[a rotação da brisa]
Coriolis gira o vetor brisa ao longo do dia: no HS,
\textbf{anti-horário}\\
noite: terral (terra esfria mais rápido) — o espelho fraco da brisa\\
hodógrafo diurno completo no Caso~1
\end{noquadro}

\subsection{Ventos de vale e encosta}

Dia: anabático (encosta aquecida puxa ar para cima) $+$ vento de vale
subindo o eixo. Noite: catabático $+$ vento de montanha descendo. O
acoplamento com a camada de inversão noturna
\javimos{5}{inversões} faz as ``piscinas frias'' de vale — as geadas
de fundo de vale do café paulista.

\section{Catabáticos}\label{sec:cat}

\begin{formadif}[o equilíbrio na encosta (T2; Caso~2)]
\[ \frac{du}{dt} = g\,\frac{\Delta\theta}{\theta}\sin\alpha
- \frac{C_D}{h}u^2
\quad\Rightarrow\quad
u_{eq} = \sqrt{\frac{g\,h\,\Delta\theta\sin\alpha}{C_D\,\theta}}. \]
Serra à noite: 2--5\,m/s. Antártida (inversão de 20\,K, camada de
300\,m, mil km de pista): 20--50\,m/s — os ventos médios mais fortes
do planeta (N2), com Coriolis desviando o escoamento encosta abaixo.
\end{formadif}

O motor é o resfriamento radiativo noturno
\javimos{2}{balanço noturno} — o mesmo do nevoeiro
\javimos{6}{nevoeiro de radiação}, agora num plano inclinado.

\section{Ondas de montanha}\label{sec:ondas}

\quadro{Retomar a mola de Brunt--Väisälä do Cap.~5: a serra é o
martelo que toca a mola. Desenhar as linhas de corrente ondulando e
as cristas inclinando contra o vento com a altura.}

\begin{formadif}[Scorer e a propagação (T3; Caso~3)]
\[ \frac{d^2\hat\delta}{dz^2} + (\ell^2 - k^2)\hat\delta = 0,
\qquad \ell^2 = \frac{N^2}{U^2} - \frac{U''}{U}. \]
Montanha larga ($k<\ell$): onda propaga para cima,
$\lambda_z = 2\pi U/N \sim 7$\,km. $\ell^2$ caindo com a altura:
ondas \emph{aprisionadas} de sotavento — o trem de lenticulares
\javimos{6}{nuvens de onda} com espaçamento $2\pi U/N$ horizontal.
\end{formadif}

\begin{noquadro}[o que as ondas fazem]
lenticulares paradas no céu com vento forte (parecem OVNIs; são
cristas)\\
rotores embaixo das cristas: perigo de aviação\\
quebra de onda em altitude: CAT \javimos{5}{Richardson} e o arrasto
de ondas que os modelos precisam parametrizar
\veremos{20}{parametrização orográfica}
\end{noquadro}

\subsection{Föhn}

Sobe úmido (chove na serra), desce seco pelo $\Gamma_d$ inteiro: chega
mais quente e seco do outro lado (T5) — já deduzido com $\theta$
\javimos{3}{föhn} e agora com o mecanismo completo: a chuva de
barlavento \emph{remove} o $L_v$ do sistema.

\section{Água rasa sobre a serra: Froude}\label{sec:hidra}

\quadro{Analogia da pia: o jato central liso (supercrítico) e o anel
turbulento (salto hidráulico). A camada fria sob inversão é a
``água''; a serra, o obstáculo.}

\begin{formadif}[regimes de Froude (T4; Caso~4)]
\[ Fr = \frac{U}{\sqrt{g'h}},\qquad g' = g\frac{\Delta\theta}{\theta}. \]
$Fr<1$: sobe engrossando (subcrítico); $Fr>1$: afina e acelera
(supercrítico); transição crítica no topo $\Rightarrow$ a camada
\textbf{despenca} supercrítica na descida e fecha num salto
hidráulico. É a física da Bora, dos \emph{windstorms} de Boulder e
das rajadas de sotavento de serra (N4).
\end{formadif}

Ventos de desfiladeiro (\emph{gap winds}): a mesma camada rasa
espremida num corredor acelera por Bernoulli — Ponta da Serra do Mar,
estreito de Gibraltar, Tehuantepec. Primo litorâneo: o \textbf{jato
costeiro de barreira} (Stull \S17.6) — ar estável bloqueado por uma
serra ($Fr < 1$!) é desviado e acelera paralelo à costa.

\section{O vento como recurso}\label{sec:eolica}

\begin{formalivro}[estatística e potência (Stull \S\S17.1--17.2)]
A frequência de velocidades num sítio segue bem a distribuição de
\textbf{Weibull},
$f(M) \propto (M/M_0)^{\alpha-1}e^{-(M/M_0)^\alpha}$; a rosa dos
ventos dá as direções. A potência disponível é
$\tfrac12\rho A M^3$ — o \emph{cubo} do vento! — e uma turbina extrai
no máximo a fração de \textbf{Betz}, $16/27 \approx 59\%$; na
prática, a curva real tem \emph{cut-in}, potência nominal e
\emph{cut-out}.
\end{formalivro}

O $M^3$ é o motivo de o perfil logarítmico decidir a altura da torre
\veremos{18}{perfil log e N2} e de os catabáticos/brisas/jatos deste
capítulo valerem dinheiro. Dentro de florestas e cidades
(\emph{canopy flows}, Stull \S17.11), o perfil log vale só acima da
copa, deslocado de uma altura $d$; dentro dela o vento decai
exponencialmente — relevante para dispersão sob dossel
\veremos{19}{dispersão}.

\section{Síntese da aula}

\begin{noquadro}[mapa final --- canto do quadro]
\[ \frac{d\Gamma}{dt} = -\oint\frac{dP}{\rho}\ \text{(brisas)} \qquad
u_{eq} = \sqrt{\tfrac{gh\Delta\theta\sin\alpha}{C_D\theta}}\
\text{(catabático)} \]
\[ \lambda_z = \frac{2\pi U}{N}\ \text{(ondas)} \qquad
Fr = \frac{U}{\sqrt{g'h}}\ \text{(vendavais de sotavento)} \]
térmicos de dia, catabáticos de noite, ondas e saltos onde há serra
\end{noquadro}

\section{Exercícios complementares}\label{sec:exercicios}

Soluções (professor) em \texttt{cap17\_solucoes.ipynb}.

\subsection*{Teóricos}

\begin{enumerate}[label=\textbf{T\arabic*.},itemsep=4pt]
  \item Enuncie o teorema de circulação de Bjerknes e estime a taxa de
        geração de circulação do circuito costa--mar para
        $\Delta T = 5$\,K numa coluna de 1\,km.
  \item Deduza a velocidade de equilíbrio catabática e avalie para uma
        encosta de serra (2°, 5\,K, 50\,m) e para a Antártida
        (0{,}6°, 20\,K, 300\,m).
  \item Apresente o parâmetro de Scorer e as condições de onda
        propagante, evanescente e aprisionada; calcule $\lambda_z$
        para $U = 12$\,m/s, $N = 0{,}011$\,s$^{-1}$.
  \item Defina o número de Froude da camada rasa e explique a
        sequência subcrítico $\to$ crítico no topo $\to$ supercrítico
        na descida $\to$ salto hidráulico.
  \item Calcule o aquecimento föhn para ar que sobe com LCL a 1\,km,
        cruza serra de 3\,km chovendo, e desce seco do outro lado
        ($T_0 = 20\,°$C).
\end{enumerate}

\subsection*{Numéricos (no notebook)}

\begin{enumerate}[label=\textbf{N\arabic*.},itemsep=4pt]
  \item Varra a forçante térmica da brisa e adicione vento sinótico
        oposto; a partir de que oposição a brisa não entra?
  \item Reproduza um catabático antártico e compare com os
        20--50\,m/s de Cape Denison.
  \item Meça $\lambda_z$ das ondas simuladas contra $2\pi U/N$ e ache
        os $U$ de ressonância na troposfera.
  \item Mapeie os regimes de Froude no plano ($Fr$, $h_m/h_0$) e
        marque a região de vendaval de sotavento.
\end{enumerate}

\vspace{1em}
\noindent\rule{\linewidth}{0.4pt}\\
{\scriptsize Material derivado de R.~Stull, \emph{Practical Meteorology}
(CC BY-NC-SA 4.0); este documento herda a mesma licença.}

\end{document}
